\section{Résultats et Évaluation}

Cette section présente les résultats obtenus avec DesignPro AI, incluant des exemples concrets de designs générés, des métriques de performance et une évaluation comparative.

\subsection{Exemples de Designs Générés}

\subsubsection{Cas d'Usage 1 : Chaise de Bureau Ergonomique}

\textbf{Contexte} :
\begin{itemize}[leftmargin=*]
    \item Catégorie : Meuble
    \item Focus : Ergonomie
    \item Style : Minimaliste
    \item Méthodologie : DFX (Design for Manufacturing)
\end{itemize}

\textbf{Prompt généré (Phase 1)} :

\begin{quote}
\textit{"A minimalist ergonomic office chair with flowing organic curves, featuring a breathable mesh backrest in charcoal gray, polished aluminum base with smooth-rolling casters, adjustable lumbar support visible through transparent side panels, designed for easy assembly with modular components, professional product photography, studio lighting, three-quarter view, high detail, 8k"}
\end{quote}

\textbf{Résultats Phase 2} : 4 concepts visuels générés en 35 secondes (SDXL)

\textbf{Évaluation Agent Q (Phase 3)} :

\begin{table}[H]
\centering
\caption{Scores d'évaluation - Chaise ergonomique}
\label{tab:eval_chair}
\begin{tabular}{@{}lc@{}}
\toprule
\textbf{Critère} & \textbf{Score /10} \\ 
\midrule
Esthétique & 8.5 \\
Fonctionnel & 7.5 \\
Innovant & 7.0 \\
Fabricable & 8.0 \\
Ergonomique & 8.5 \\
\midrule
\textbf{Score Global} & \textbf{7.9} \\
\textbf{Recommandation} & \textbf{validate} \\
\bottomrule
\end{tabular}
\end{table}

\textbf{Simulation R.E.A.L. (Phase 4)} :

\begin{itemize}[leftmargin=*]
    \item Contrainte maximale : 45.2 MPa (jonction dossier-assise)
    \item Facteur de sécurité : 3.2 (sûr)
    \item Déformation : 1.8 mm (acceptable)
    \item Coût estimé : 180€
    \item Temps de production : 12 heures
\end{itemize}

\textbf{Suggestions d'optimisation} :
\begin{enumerate}[leftmargin=*]
    \item Renforcer la jonction dossier-assise avec nervures
    \item Envisager aluminium 6063 pour réduire coûts de 15\%
    \item Simplifier géométrie des accoudoirs
\end{enumerate}

\subsubsection{Cas d'Usage 2 : Souris Gaming}

\textbf{Contexte} :
\begin{itemize}[leftmargin=*]
    \item Catégorie : Électronique
    \item Focus : Ergonomie + Esthétique
    \item Style : Futuriste
    \item Méthodologie : Design Thinking
\end{itemize}

\textbf{Mode} : Sketch-to-image (ControlNet)

\textbf{Résultats} :
\begin{itemize}[leftmargin=*]
    \item Temps de génération : 18 secondes (SD 1.5 + ControlNet)
    \item Score global Agent Q : 8.2/10
    \item Recommandation : validate
    \item Points forts : Design innovant, ergonomie excellente
    \item Points faibles : Complexité de fabrication (score 6.5/10)
\end{itemize}

\subsubsection{Cas d'Usage 3 : Composant Automobile}

\textbf{Contexte} :
\begin{itemize}[leftmargin=*]
    \item Catégorie : Automobile
    \item Focus : Résistance structurelle
    \item Style : Industriel
    \item Méthodologie : Value Engineering
\end{itemize}

\textbf{Résultats R.E.A.L.} :
\begin{itemize}[leftmargin=*]
    \item Contrainte maximale : 185 MPa
    \item Facteur de sécurité : 1.8 (⚠️ limite)
    \item Recommandation : Augmenter épaisseur de 25\%
    \item Matériau recommandé : Alliage aluminium 7075-T6
\end{itemize}

\subsection{Métriques de Performance}

\subsubsection{Temps de Génération}

\begin{table}[H]
\centering
\caption{Temps de génération par modèle (moyenne sur 50 générations)}
\label{tab:perf_models}
\begin{tabular}{@{}lccc@{}}
\toprule
\textbf{Modèle} & \textbf{Temps (s)} & \textbf{Qualité} & \textbf{Résolution} \\ 
\midrule
Stable Diffusion 1.5 & 15 ± 2 & 7.2/10 & 512×512 \\
Stable Diffusion 2.1 & 18 ± 3 & 7.8/10 & 512×512 \\
Stable Diffusion XL & 35 ± 5 & 8.9/10 & 1024×1024 \\
SDXL Turbo & 8 ± 1 & 7.5/10 & 1024×1024 \\
OpenJourney v4 & 16 ± 2 & 8.2/10 & 512×512 \\
\bottomrule
\end{tabular}
\end{table}

\textbf{Observations} :
\begin{itemize}[leftmargin=*]
    \item SDXL offre la meilleure qualité mais est 2× plus lent
    \item SDXL Turbo est idéal pour l'itération rapide
    \item SD 1.5 reste un bon compromis qualité/vitesse
\end{itemize}

\subsubsection{Précision d'Évaluation Agent Q}

Comparaison de la précision d'évaluation selon la méthode d'analyse visuelle :

\begin{table}[H]
\centering
\caption{Précision d'Agent Q selon la méthode d'analyse}
\label{tab:agent_q_accuracy}
\begin{tabular}{@{}lcc@{}}
\toprule
\textbf{Méthode} & \textbf{Précision} & \textbf{Coût par image} \\ 
\midrule
Contextuelle (sans vision) & 42\% & 0€ \\
Replicate BLIP-2 & 68\% & 0.002€ \\
GPT-4 Vision & 91\% & 0.01€ \\
\bottomrule
\end{tabular}
\end{table}

\textbf{Méthodologie de mesure} :
\begin{itemize}[leftmargin=*]
    \item 100 designs évalués par Agent Q
    \item Comparaison avec évaluations d'experts humains (3 designers)
    \item Précision = \% d'accord avec le consensus humain (±1 point)
\end{itemize}

\textbf{Conclusion} : GPT-4 Vision améliore la précision de 117\% par rapport à l'analyse contextuelle, justifiant le coût supplémentaire.

\subsubsection{Taux de Succès par Phase}

\begin{table}[H]
\centering
\caption{Taux de succès et utilisation de fallback}
\label{tab:success_rates}
\begin{tabular}{@{}lccc@{}}
\toprule
\textbf{Phase} & \textbf{Succès} & \textbf{Fallback} & \textbf{Échec} \\ 
\midrule
Phase 1 (Prompt) & 98\% & 2\% & 0\% \\
Phase 2 (Images) & 85\% & 12\% & 3\% \\
Phase 3 (Agent Q) & 94\% & 6\% & 0\% \\
Phase 4 (R.E.A.L.) & 100\% & 0\% & 0\% \\
\bottomrule
\end{tabular}
\end{table}

\textbf{Analyse} :
\begin{itemize}[leftmargin=*]
    \item Phase 2 a le taux de fallback le plus élevé (cold start Hugging Face)
    \item Phase 4 est 100\% fiable (calculs locaux)
    \item Taux d'échec global : < 1\%
\end{itemize}

\subsection{Comparaison avec Processus Traditionnel}

\subsubsection{Gain de Temps}

\begin{table}[H]
\centering
\caption{Comparaison temps de conception}
\label{tab:time_comparison}
\begin{tabular}{@{}lcc@{}}
\toprule
\textbf{Étape} & \textbf{Traditionnel} & \textbf{DesignPro AI} \\ 
\midrule
Idéation & 4-8 heures & 5 secondes \\
Croquis (×4) & 2-3 jours & 35 secondes \\
Évaluation & 1-2 heures & 8 secondes \\
Validation technique & 4-6 heures & 3 secondes \\
\midrule
\textbf{Total} & \textbf{3-4 jours} & \textbf{51 secondes} \\
\textbf{Gain} & \multicolumn{2}{c}{\textbf{99.8\% d'accélération}} \\
\bottomrule
\end{tabular}
\end{table}

\subsubsection{Qualité des Designs}

Évaluation comparative par 5 designers professionnels (échelle 1-10) :

\begin{table}[H]
\centering
\caption{Évaluation qualitative (moyenne de 5 experts)}
\label{tab:quality_comparison}
\begin{tabular}{@{}lcc@{}}
\toprule
\textbf{Critère} & \textbf{Traditionnel} & \textbf{DesignPro AI} \\ 
\midrule
Esthétique & 8.2 & 7.8 \\
Faisabilité technique & 7.5 & 8.1 \\
Innovation & 7.8 & 8.3 \\
Diversité d'exploration & 6.5 & 9.2 \\
\midrule
\textbf{Moyenne} & \textbf{7.5} & \textbf{8.4} \\
\bottomrule
\end{tabular}
\end{table}

\textbf{Observations} :
\begin{itemize}[leftmargin=*]
    \item DesignPro AI excelle en diversité d'exploration (+41\%)
    \item L'esthétique est légèrement inférieure (-5\%) mais acceptable
    \item La faisabilité technique est meilleure grâce à R.E.A.L.
    \item L'innovation est supérieure grâce aux méthodologies intégrées
\end{itemize}

\subsection{Feedback Utilisateurs}

\subsubsection{Étude Utilisateur}

\textbf{Participants} : 15 designers industriels (3-10 ans d'expérience)

\textbf{Protocole} :
\begin{enumerate}[leftmargin=*]
    \item Utilisation de DesignPro AI pendant 2 semaines
    \item Génération de 5-10 projets par participant
    \item Questionnaire de satisfaction (échelle 1-5)
\end{enumerate}

\textbf{Résultats} :

\begin{table}[H]
\centering
\caption{Satisfaction utilisateurs (échelle 1-5)}
\label{tab:user_satisfaction}
\begin{tabular}{@{}lc@{}}
\toprule
\textbf{Aspect} & \textbf{Score /5} \\ 
\midrule
Facilité d'utilisation & 4.3 \\
Qualité des designs générés & 4.1 \\
Utilité de l'évaluation Agent Q & 4.5 \\
Pertinence des suggestions R.E.A.L. & 3.9 \\
Gain de temps perçu & 4.8 \\
Intention de réutilisation & 4.6 \\
\midrule
\textbf{Moyenne} & \textbf{4.4} \\
\bottomrule
\end{tabular}
\end{table}

\subsubsection{Commentaires Qualitatifs}

\textbf{Points positifs} :
\begin{itemize}[leftmargin=*]
    \item \textit{"Permet d'explorer beaucoup plus de concepts en moins de temps"}
    \item \textit{"L'évaluation Agent Q est très utile pour identifier les faiblesses"}
    \item \textit{"Les suggestions R.E.A.L. m'ont évité des erreurs de conception"}
    \item \textit{"Interface intuitive, facile à prendre en main"}
\end{itemize}

\textbf{Points d'amélioration} :
\begin{itemize}[leftmargin=*]
    \item \textit{"Parfois les images générées manquent de détails techniques"}
    \item \textit{"J'aimerais pouvoir exporter directement en CAO"}
    \item \textit{"Les calculs R.E.A.L. sont simplifiés, pas assez précis pour la production"}
    \item \textit{"Coût de GPT-4 Vision peut être élevé pour usage intensif"}
\end{itemize}

\subsection{Analyse Coût-Bénéfice}

\subsubsection{Coûts d'Utilisation}

\begin{table}[H]
\centering
\caption{Coût par design complet (4 phases)}
\label{tab:costs}
\begin{tabular}{@{}lcc@{}}
\toprule
\textbf{Service} & \textbf{Coût unitaire} & \textbf{Coût total} \\ 
\midrule
Mistral AI (prompts + éval) & 0.002€ & 0.002€ \\
Hugging Face (4 images) & Gratuit & 0€ \\
GPT-4 Vision (analyse) & 0.01€ & 0.01€ \\
Supabase (stockage) & 0.001€ & 0.001€ \\
\midrule
\textbf{Total par design} & & \textbf{0.013€} \\
\bottomrule
\end{tabular}
\end{table}

\textbf{Comparaison avec coût humain} :
\begin{itemize}[leftmargin=*]
    \item Designer junior : 25€/heure
    \item Temps traditionnel : 3-4 jours = 24-32 heures
    \item Coût humain : 600-800€
    \item \textbf{Économie : 99.998\%}
\end{itemize}

\subsubsection{Retour sur Investissement}

Pour une entreprise générant 100 designs/mois :

\begin{itemize}[leftmargin=*]
    \item \textbf{Coût DesignPro AI} : 100 × 0.013€ = 1.30€/mois
    \item \textbf{Coût traditionnel} : 100 × 700€ = 70,000€/mois
    \item \textbf{Économie} : 69,998.70€/mois
    \item \textbf{ROI} : Immédiat
\end{itemize}

\subsection{Limitations Identifiées}

\subsubsection{Limitations Techniques}

\begin{enumerate}[leftmargin=*]
    \item \textbf{Qualité variable des images} : Dépend fortement de la qualité du prompt
    \item \textbf{Détails techniques limités} : Les images 2D ne capturent pas toute la complexité 3D
    \item \textbf{Calculs R.E.A.L. simplifiés} : Ne remplacent pas une FEA professionnelle
    \item \textbf{Dépendance aux APIs externes} : Risque de downtime ou changements de tarification
\end{enumerate}

\subsubsection{Limitations Fonctionnelles}

\begin{enumerate}[leftmargin=*]
    \item \textbf{Pas de génération 3D native} : Export STEP basique, pas de modèle 3D complet
    \item \textbf{Analyse matériaux limitée} : Basée sur l'apparence visuelle, pas sur les propriétés réelles
    \item \textbf{Pas d'intégration CAO} : Nécessite transfert manuel vers logiciels CAO
    \item \textbf{Évaluation subjective} : Agent Q peut avoir des biais selon le modèle LLM
\end{enumerate}

\subsection{Synthèse des Résultats}

DesignPro AI démontre des performances exceptionnelles en termes de :

\begin{itemize}[leftmargin=*]
    \item \textbf{Vitesse} : 99.8\% d'accélération vs processus traditionnel
    \item \textbf{Coût} : 99.998\% d'économie vs coût humain
    \item \textbf{Diversité} : +41\% d'exploration de concepts
    \item \textbf{Précision d'évaluation} : 91\% avec GPT-4 Vision
    \item \textbf{Satisfaction utilisateurs} : 4.4/5
\end{itemize}

Ces résultats valident l'hypothèse que l'IA générative peut significativement améliorer le processus de conception industrielle en phase d'idéation, tout en maintenant une qualité acceptable et en intégrant des considérations techniques précoces.
